\documentclass[12pt,letterpaper]{article}

\usepackage[T1]{fontenc}
\usepackage[utf8]{inputenc}
\usepackage[spanish,mexico,es-tabla,es-noshorthands]{babel}
\usepackage{newtxtext,newtxmath}          % Times New Roman
\usepackage[margin=2.5cm]{geometry}
\usepackage{graphicx}
\usepackage{float}
\usepackage{caption}
\usepackage{subcaption}
\usepackage{booktabs}
\usepackage{enumitem}
\usepackage{amsmath}          % no cargar amssymb: choca con \Bbbk de newtxmath
\usepackage{listings}
\usepackage{xcolor}
\usepackage{setspace}
\usepackage[colorlinks=true,allcolors=black]{hyperref}
\usepackage{todonotes}                    % usar [disable] para la versión final

\graphicspath{{evidencias/}{../assets/}}
\captionsetup{labelsep=colon,font=small,labelfont=bf}
\onehalfspacing
\setlength{\parskip}{6pt}

% Títulos de sección al estilo del formato de la profesora
\usepackage{titlesec}
\titleformat{\section}{\normalfont\large\bfseries}{}{0pt}{}
\titleformat{\subsection}{\normalfont\normalsize\bfseries}{}{0pt}{}

\definecolor{codekw}{RGB}{0,0,180}
\definecolor{codecm}{RGB}{0,128,0}
\definecolor{codestr}{RGB}{170,0,0}
\definecolor{codebg}{RGB}{248,248,248}

\lstdefinestyle{hdl}{
  language=VHDL,
  basicstyle=\ttfamily\footnotesize,
  keywordstyle=\color{codekw}\bfseries,
  commentstyle=\color{codecm}\itshape,
  stringstyle=\color{codestr},
  backgroundcolor=\color{codebg},
  numbers=left, numberstyle=\tiny\color{gray}, numbersep=8pt,
  frame=single, rulecolor=\color{gray!50},
  breaklines=true, breakatwhitespace=true,
  tabsize=2, showstringspaces=false, captionpos=b,
  extendedchars=true, inputencoding=utf8,
  literate={á}{{\'a}}1 {é}{{\'e}}1 {í}{{\'i}}1 {ó}{{\'o}}1 {ú}{{\'u}}1
           {Á}{{\'A}}1 {É}{{\'E}}1 {Í}{{\'I}}1 {Ó}{{\'O}}1 {Ú}{{\'U}}1
           {ñ}{{\~n}}1 {Ñ}{{\~N}}1 {ü}{{\"u}}1 {¿}{{?`}}1 {¡}{{!`}}1
}
\lstset{style=hdl}
% Para que el código comparta numeración con las figuras (en vez de tener su propio
% contador "Listing"), NO usar caption=/label= dentro de lstlisting/lstinputlisting.
% En su lugar, envolver el listado en un \begin{figure} normal y poner \caption y
% \label ahí (\figurename ya es "Figura" por babel-spanish):
%   \begin{figure}[H]
%     \centering
%     \lstinputlisting[style=hdl]{codigo/archivo.vhd}
%     \caption{Descripción del código.}
%     \label{fig:archivo}
%   \end{figure}


\usepackage{colortbl}
\definecolor{kgroupa}{RGB}{190,220,255}
\definecolor{kgroupb}{RGB}{255,224,178}

\begin{document}

\begin{titlepage}
\centering
\noindent
\begin{minipage}[c]{0.45\textwidth}\centering
  \includegraphics[height=3.2cm]{logo-unam.jpg}
\end{minipage}\hfill
\begin{minipage}[c]{0.45\textwidth}\centering
  \includegraphics[height=3.2cm]{logo-fi.png}
\end{minipage}

\vspace{1.2cm}
{\Large\bfseries UNIVERSIDAD NACIONAL AUTÓNOMA\\[2pt] DE MÉXICO\par}
\vspace{0.9cm}
{\large\bfseries FACULTAD DE INGENIERÍA\par}
\vspace{0.6cm}
{\large\bfseries INGENIERÍA EN COMPUTACIÓN\par}
\vspace{0.6cm}
{\bfseries Laboratorio de Organización de Arquitectura de Computadoras\par}
\vspace{0.9cm}
{\bfseries Reporte - Construcción de máquinas de Estados usando memorias: Direccionamiento por Trayectoria\par}

\vspace{1.0cm}
\raggedright
{\bfseries NOMBRE DE LOS INTEGRANTES:}\\
Omar Alfonso Zendejas Labias \quad Grupo de Teoría: 04\\
Emiliano García Olvera \quad Grupo de Teoría: 02

\vspace{1.0cm}
{\bfseries NOMBRE DEL PROFESOR: ING. KAREN PÉREZ}

\vspace{0.6cm}
{\bfseries GRUPO: 06 \quad EQUIPO: 9}

\vspace{0.6cm}
{\bfseries SEMESTRE:} \underline{2027-1.}

\vspace{0.6cm}
\hfill {\bfseries FECHA DE ENTREGA: 08 de octubre de 2026}
\thispagestyle{empty}
\end{titlepage}

\setcounter{page}{1}

\section{Objetivo}

Familiarizar al alumno con la construcción de máquinas de estados mediante el uso de
memorias, aplicando el método de direccionamiento por trayectoria.

\end{document}